\chapter{روش‌شناسی}
\label{chap:method}

\section{جریان کلی آزمایش}
شکل~\ref{fig:experiment-flow} جریان آزمایش را نشان می‌دهد. ابتدا از سری بار و \lr{PV}، برنامه مرجع مقید با برنامه‌ریزی پویا ساخته می‌شود. سپس دو شبکه با این داده آموزش می‌بینند. تفاوت آن دو در نحوه تعریف خروجی است: شبکه داده‌محور چهار خروجی آزاد می‌دهد، اما شبکه فیزیک‌تعبیه‌شده فقط فرمان باتری را می‌سازد و همه متغیرهای دیگر از قوانین فصل قبل حاصل می‌شوند. در انتها، هر دو مدل روی افق آزمون با \lr{SOC} تولیدشده توسط خودشان بازگشت داده می‌شوند.

\begin{figure}[H]
  \centering
  \begin{tikzpicture}[node distance=0.85cm, every node/.style={align=center}]
    \node[block, minimum width=5.4cm] (data) {داده عمومی ترکیبی\\بار، \lr{PV}، تعرفه و زمان};
    \node[block, below=of data, minimum width=5.4cm] (dp) {برنامه‌ریزی پویا\\برنامه مرجع مقید};
    \node[block, below=of dp, minimum width=5.4cm] (nets) {آموزش دو شبکه\\داده‌محور و فیزیک‌تعبیه‌شده};
    \node[loss, below=of nets, minimum width=5.4cm] (layer) {لایه فیزیکی قطعی\\توان، \lr{SOC} و توازن};
    \node[loss, below=of layer, minimum width=5.4cm, fill=yellow!8] (eval) {آزمون بازگشتی\\امکان‌پذیری و هزینه};
    \draw[flow] (data) -- (dp);
    \draw[flow] (dp) -- (nets);
    \draw[flow] (nets) -- (layer);
    \draw[flow] (layer) -- (eval);
  \end{tikzpicture}
  \caption{جریان مطالعه موردی و مقایسه کنترل‌کننده‌ها. لایه فیزیکی در مدل فیزیک‌تعبیه‌شده بخشی از خودِ خروجی است؛ برای مدل داده‌محور فقط در زنجیره ترمیم‌شده هزینه اعمال می‌شود.}
  \label{fig:experiment-flow}
\end{figure}

\section{بردار ورودی و مرجع آموزشی}
برای ساعت $t$، بردار ورودی شامل زمان تناوبی، وضعیت روز، بار، \lr{PV}، تعرفه، \lr{SOC} و پنجره پیش‌بینی کوتاه است:
\begin{equation}
\begin{split}
  \bm{x}(t)=\bigl(&\sin(2\pi h/24),\,\cos(2\pi h/24),\,d,\,P_{\mathrm{L}},\,P_{\mathrm{PV}},\,c_{\mathrm{b}},\,c_{\mathrm{s}},\,s,\\
  &P_{\mathrm{L}}^{t+1},\,P_{\mathrm{PV}}^{t+1},\,P_{\mathrm{L}}^{t+3},\,P_{\mathrm{PV}}^{t+3}\bigr).
\end{split}
  \label{eq:model-input}
\end{equation}
$h$ ساعت روز و $d$ نشانگر روز کاری است. مقادیر یک و سه ساعت آینده، پیش‌بینی بی‌خطا فرض می‌شوند. این فرض برای مقایسه آفلاین با برنامه مرجع است و ادعای دسترسی بی‌خطا به پیش‌بینی در میدان نیست.

برنامه‌ریزی پویا با مشاهده کل افق برون‌زاد، دنباله مرجع
\begin{equation}
  \bm{y}^{\star}(t)=\bigl(P_{\mathrm{G}}^{\star},\,P_{\mathrm{c}}^{\star},\,P_{\mathrm{d}}^{\star},\,s^{\star}(t+\Delta t)\bigr)
  \label{eq:reference-output}
\end{equation}
را می‌سازد. آماره‌های استانداردسازی، یعنی میانگین و انحراف معیار، فقط از مجموعه آموزش محاسبه می‌شوند تا از نشت آماری مجموعه اعتبارسنجی و آزمون جلوگیری شود.

\section{شبکه داده‌محور آزاد}
مدل نخست یک شبکه پیش‌خور با دو لایه پنهان ۳۲ واحدی و فعال‌ساز تانژانت هیپربولیک است. خروجی آن چهار مؤلفه~\eqref{eq:reference-output} را مستقیماً پیش‌بینی می‌کند. تابع آموزش آن میانگین مربع خطای خروجی استانداردشده است. این مدل برای آن در نظر گرفته شده است که نشان دهد نزدیک بودن به برچسب، به‌تنهایی معادلِ قابل‌اجرا بودن برنامه نیست.

خروجی خام برای سنجش نقض‌ها با \lr{SOC} پیش‌بینی‌شده خود مدل بازگشت داده می‌شود. در زنجیره اجرایی ترمیم‌شده، \lr{SOC} ورودی هر گام از فرمان نگاشت‌یافته و دینامیک باتری دوباره محاسبه می‌شود، نه از مؤلفه خام \lr{SOC} شبکه. برای محاسبه هزینه منصفانه، فرمان خام مدل داده‌محور با نگاشت ایمنی پس‌پردازشی
\begin{equation}
  u_{\mathrm{safe}}=\mathrm{clip}(\hat{P}_{\mathrm{c}}-\hat{P}_{\mathrm{d}},-P_{\max},P_{\max})
  \label{eq:usafe}
\end{equation}
به لایه~\eqref{eq:hard-soc} داده می‌شود. عملگر گیره به‌صورت $\mathrm{clip}(x,a,b)=\min(\max(x,a),b)$ تعریف می‌شود. هزینه حاصل با برچسب «داده‌محور با نگاشت ایمنی» گزارش می‌شود و نباید با هزینه خروجی خامِ احتمالاً ناممکن اشتباه گرفته شود.

\section{شبکه فیزیک‌تعبیه‌شده}
شبکه دوم همان عرض و عمق لایه‌های پنهان را دارد، ولی به‌جای چهار خروجی آزاد، فقط یک عدد $z(t)$ می‌سازد. فرمان علامت‌دار باتری از رابطه
\begin{equation}
  u(t)=P_{\max}\tanh\bigl(z(t)\bigr)
  \label{eq:signed-command}
\end{equation}
به‌دست می‌آید؛ علامت مثبت برای شارژ و منفی برای دشارژ است. چون بازه $\tanh$ در $\pm 1$ باز است، فرمان دقیقاً برابر $\pm P_{\max}$ فقط به‌صورت حدی نزدیک می‌شود. سقف وابسته به \lr{SOC} برای شارژ و دشارژ چنین تعریف می‌شود:
\begin{equation}
\begin{aligned}
  \bar{P}_{\mathrm{c}}(s)&=\max\left(0,\min\left(P_{\max},\frac{(s_{\max}-s)E_{\max}}{\eta_{\mathrm{c}}\Delta t}\right)\right),\\
  \bar{P}_{\mathrm{d}}(s)&=\max\left(0,\min\left(P_{\max},\frac{(s-s_{\min})E_{\max}\eta_{\mathrm{d}}}{\Delta t}\right)\right).
\end{aligned}
  \label{eq:soc-caps}
\end{equation}
سپس خروجی‌ها به‌صورت
\begin{equation}
  P_{\mathrm{c}}=\min\bigl([u]_{+},\bar{P}_{\mathrm{c}}(s)\bigr),\qquad
  P_{\mathrm{d}}=\min\bigl([-u]_{+},\bar{P}_{\mathrm{d}}(s)\bigr)
  \label{eq:hard-powers}
\end{equation}
و
\begin{equation}
\begin{aligned}
  s(t+\Delta t)&=s(t)+\frac{\eta_{\mathrm{c}}P_{\mathrm{c}}\Delta t}{E_{\max}}-\frac{P_{\mathrm{d}}\Delta t}{\eta_{\mathrm{d}}E_{\max}},\\
  P_{\mathrm{G}}&=P_{\mathrm{L}}+P_{\mathrm{c}}-P_{\mathrm{PV}}-P_{\mathrm{d}}
\end{aligned}
  \label{eq:hard-soc}
\end{equation}
ساخته می‌شوند. در این مطالعه برای $P_{\mathrm{G}}$ سقف واردات یا صادرات در نظر گرفته نشده است؛ این انتخاب مدل‌سازی باعث می‌شود توازن توان پس از تعیین فرمان باتری همواره با تبادل شبکه بسته شود.

از~\eqref{eq:hard-powers} روشن است که در هر گام حداکثر یکی از $P_{\mathrm{c}}$ و $P_{\mathrm{d}}$ مثبت است؛ سقف‌های~\eqref{eq:soc-caps} با $\max(0,\cdot)$ نامنفی می‌مانند حتی اگر $s$ به‌خاطر گرد کردن اندکی از بازه خارج شود. از~\eqref{eq:soc-caps} و~\eqref{eq:hard-soc} نیز حدود \lr{SOC} و توازن توان برقرار می‌شوند. بنابراین باقیمانده‌های~\eqref{eq:power-residual} و دینامیک \lr{SOC} در این مدل باید، به جز خطای گرد کردن، صفر باشند. این تضمین مشروط به آن است که وضعیت شارژ آغازین در $[s_{\min},s_{\max}]$ باشد و در آزمون بازگشتی نیز $s$ از همین نگاشت به‌روز شود. تابع آموزش این شبکه خطای فرمان مرجع است:
\begin{equation}
  \mathcal{L}_{u}=\frac{1}{N}\sum_{i=1}^{N}\left(\tanh(z_{i})-\frac{u_{i}^{\star}}{P_{\max}}\right)^{2}.
  \label{eq:command-loss}
\end{equation}
این ساختار از نظر مفهومی با یادگیری فیزیک‌آگاه هم‌خوان است، اما اصطلاح دقیق آن در این گزارش «فیزیک‌تعبیه‌شده با قید سخت» است. مزیت آن نسبت به جریمه نرم این است که وزن جریمه برای تضمین قیود اساسی لازم نیست. ترتیب محاسبات لایه در پیوست~\ref{app:hard-layer} آمده است.

\section{تنظیم آموزش و افراز داده}
نمونه‌ها به ترتیب زمانی حفظ می‌شوند: ۶۰ درصد اول برای آموزش، ۲۰ درصد بعدی برای اعتبارسنجی و ۲۰ درصد پایانی برای آزمون است. شبکه‌ها با الگوریتم آدام، دسته‌های ۶۴تایی و حداکثر ۴۵۰ \lr{epoch} آموزش می‌بینند. توقف زودهنگام با صبر ۵۵ \lr{epoch} فقط بر خطای اعتبارسنجی اعمال می‌شود. پنج بذر ۷، ۱۱، ۱۹، ۲۳ و ۲۹ فقط مقداردهی اولیه و ترتیب دسته‌ها را تغییر می‌دهند. در نتیجه، انحراف معیار جدول نتایج پراکندگی آموزش را نشان می‌دهد، نه تغییر خود داده یا تغییر تقسیم داده. در آموزش، مؤلفه $s$ از مسیر مرجع است؛ در آزمون بازگشتی از مسیر خود مدل می‌آید. بنابراین بخشی از شکاف هزینه نسبت به مرجع می‌تواند ناشی از این ناهماهنگی توزیع باشد، نه فقط از خطای تقلید یک‌گامی.

\section{شاخص‌های ارزیابی}
برای شفافیت، \lr{MSE} تجمیعی در این گزارش صریحاً میانگین روی چهار مؤلفه خروجی است:
\begin{equation}
  \mathrm{MSE}_{\mathrm{comp}}=\frac{1}{4N}\sum_{i=1}^{N}\sum_{j=1}^{4}\bigl(\hat{y}_{ij}-y_{ij}^{\star}\bigr)^{2}.
  \label{eq:mse-comp}
\end{equation}
از آن‌جا که توان و \lr{SOC} هم‌واحد نیستند، این عدد به‌تنهایی معیار فیزیکی کافی نیست. به همین دلیل، میانگین قدر مطلق خطای توان شبکه، فرمان باتری و \lr{SOC} نیز جداگانه گزارش می‌شوند. برای مدل فیزیک‌تعبیه‌شده، چهار مؤلفه پس از لایه سخت ساخته می‌شوند و تابع آموزش همان~\eqref{eq:command-loss} است، نه این میانگین چهارتایی.

امکان‌پذیری با پنج شاخص کنترل می‌شود: میانگین $|r_{\mathrm{G}}|$، نرخ نقض \lr{SOC}، نرخ توان شارژ منفی، نرخ توان دشارژ منفی و نرخ شارژ و دشارژ هم‌زمان. برای هر برنامه قابل اجرا، هزینه کل آزمون نیز از جمع~\eqref{eq:step-cost} محاسبه می‌گردد. مرجع برنامه‌ریزی پویا و حالت بدون باتری، دو مرز قابل تفسیر برای هزینه هستند.
